\textbf{PHP} is one of the most widely used server-side scripting languages, powering approximately 77\% of all websites globally, including platforms like WordPress, Drupal, and Joomla. Its popularity stems from its ease of use, extensive library support, and seamless integration with web servers and databases, making it a top choice for web development. 
PHP language support various syntaxs which enables highly flexible and adaptable coding patterns. For example, variable variable and variable functions allow developers to access a variable and call functions based on dynamically determined names at runtime: \lstinline[language=myPHP]| \${\$varName}; \$funcName();|. 
Runtime file inclusions dynamically include PHP scripts based on variables or user inputs, enabling modularity but introducing potential risks: \lstinline[language=PHP]|include $filePath;|. 
In addition, object-oriented programming (OOP) mechanisms introduce dynamic behavior through polymorphism and dynamic method resolution: \lstinline[language=myPHP]|new ClsName() -> meth()|. 
It also supports constructing and executing code dynamically using mechanisms like the eval function or reflection APIs: \lstinline[language=myPHP]|eval('$result = $a + $b;');|. 
However, such strong and fixable features also lead to the security issue. 

\textbf{Taint-style vulnerability} occurs when untrusted user input flows through a program and reaches sensitive operations without proper validation or sanitization. These vulnerabilities are a significant concern in PHP web applications due to reliance on user inputs, which can originate from untrusted sources. If such inputs are used directly in operations like SQL queries, HTML output, or file paths, they can enable attacks such as SQL injection, cross-site scripting (XSS), or remote file inclusion (RFI). 

\textbf{Static Taint Analysis} is a program analysis technique used to identify taint-type vulnerabilities by tracing the flow of untrusted data through a program without executing it. It defines \textbf{sources}, where untrusted data enters the program (e.g., \lstinline[language=myPHP]| $\_GET, $\_POST|), and \textbf{sinks}, which are sensitive operations that may be exploited if such data is not properly handled (e.g., database queries, file operations, or system commands). The analysis tracks how tainted data propagates through variables, function calls, and control flow branches. When tainted data reaches a sink without passing through validation or sanitization, it is flagged as a potential vulnerability. 
However, taint analysis struggles to handle PHP’s dynamic and flexible features because these syntaxes introduce behaviors that depend on runtime data, which static tools often cannot accurately resolve. For instance, in the case of dynamic function calls, where the function name is determined at runtime, tools must infer the value of from its context. Without this runtime information, the target function cannot be identified, breaking the taint propagation. Similarly, variable variables, dynamically determine variable names, creating dependencies that tools must track through multiple assignments and control flows. This increases the complexity of analyzing data propagation, as traditional static tools rely on fixed variable names and structures. 

\textbf{Abstract Interpretation} is a static analysis method that approximates program execution by mapping concrete program states, such as variable values, to abstract representations that capture relevant properties. 
Instead of evaluating the exact behavior of a program, abstract interpretation operates over an abstract domain, which simplifies the program’s state space. For example, instead of tracking specific integers, it might classify them as positive, negative, or zero. During analysis, the tool explores the program's control flow, propagates abstract states through its operations, and computes a stable state, known as a fixpoint, to represent loops or recursive behavior. By systematically analyzing  possible execution paths, abstract interpretation determines how data flows through the program, identifying potential vulnerabilities like user input reaching sensitive operations. 
Taint analysis can be built on top of abstract interpretation by using abstract domains to model taint propagation. Abstract interpretation provides the underlying mechanism to track data flow precisely across complex control flow structures, making it a powerful tool for implementing robust taint analysis in dynamic languages like PHP. This approach is particularly effective for PHP because it models dynamic features such as variable function names or dynamic includes without needing to execute the program, enabling the detection of security risks in a scalable and precise manner.

\textbf{Challenge of Abstract Interpretation on PHP}. There are two possible challenges to implement the abstract interpretation. Firstly, how to precisely model the semantics of PHP syntax, particularly for dynamic features like function calls, file inclusions, and object-oriented constructs. Incomplete or imprecise modeling of these constructs may lead to false positives (erroneously identifying issues) or false negatives (failing to detect actual vulnerabilities). Furthermore, PHP support ton of syntax incrementally in over ten years updates, so it demands substantial engineering effort and theoretical advancements to support those syntax.
Secondly, scalability poses another challenge for abstract interpretation, especially when applied to large-scale or complex systems. The design of abstraction domains plays a key role in determining the trade-offs between precision and efficiency. Fine-grained abstractions improve accuracy but incur higher computational costs, while coarse-grained abstractions sacrifice precision for performance, potentially missing critical details. Additionally, analyzing intricate control flows, deep nesting, and recursive structures often leads to exponential growth in computational overhead. For instance, resolving multiple levels of class inheritance or tracking object property dependencies across deep function call stacks can be computationally expensive. These challenges limit the practical application of abstract interpretation in scenarios requiring real-time analysis or handling extensive codebases, making scalability a pivotal consideration in its development.


\textbf{Overall idea on our approach}. In this paper, our tool is a specialized static analysis tool designed to detect vulnerabilities in PHP applications by leveraging abstract interpretation and graph-based modeling. It systematically processes source code through preprocessing, parsing, and two phases of interpretation. During analysis, it constructs Object Dependence Graphs (ODGs) and Code Property Graphs (CPGs) to model both intra- and inter-procedural data and control flows. 

By abstracting PHP’s dynamic features, such as variable function names and dynamic includes, ODGen-PHP accurately tracks tainted data from sources to sinks. Its integration of backward slicing reduces computational overhead by focusing only on critical control flows. This approach enables efficient and precise detection of vulnerabilities like SQL injection and command injection, making it a robust solution for securing complex PHP codebases.